% Source: https://artofproblemsolving.com/community/c3h1207406_altitudes
% Stable identity: aops-topic-1207406-q01
By Heron's formula with $s=9$, the area is $\sqrt{9\cdot4\cdot3\cdot2}=\sqrt{216}=6\sqrt6$. The longest altitude is the one to the shortest side, $5$:

\begin{asy}
usepackage("amsmath");
size(3.5cm);
defaultpen(fontsize(10pt));
pair A=(0,0),B=(5,0); pair C=(1.2,sqrt(36-1.44)); draw(A--B--C--cycle,linewidth(0.9)); draw(C--(C.x,0),blue);
label("$5$",(2.5,0),S); label("$6$",(A+C)/2,W); label("$7$",(B+C)/2,NE); label("$h$",(C.x,2.5),E);
label("\mbox{}",(min(currentpicture,user=true)+max(currentpicture,user=true))/2,black); // keeps the text after the figure black when LaTeX typesets the labels inline
\end{asy}

\[h=\dfrac{2\cdot6\sqrt6}{5}=\boxed{\dfrac{12\sqrt6}{5}}\approx5.88.\]
